\documentclass[11pt, a4paper]{article}

% ── Page layout ───────────────────────────────────────────────────────────────
\usepackage[top=1in, bottom=1in, left=1.1in, right=1.1in]{geometry}

% ── Typography ────────────────────────────────────────────────────────────────
\usepackage{mathptmx}          % Times New Roman font
\usepackage{microtype}         % Better justification
\usepackage{parskip}           % Space between paragraphs instead of indent
\setlength{\parskip}{6pt}

% ── Mathematics ───────────────────────────────────────────────────────────────
\usepackage{amsmath}
\usepackage{amssymb}

% ── Figures and tables ────────────────────────────────────────────────────────
\usepackage{graphicx}
\graphicspath{figures/}
\usepackage{booktabs}
\usepackage{caption}
\captionsetup{font=small, labelfont=bf}

% ── Section formatting ────────────────────────────────────────────────────────
\usepackage{titlesec}
\titleformat{\section}{\large\bfseries}{\thesection.}{0.6em}{}
\titleformat{\subsection}{\normalsize\bfseries}{\thesubsection.}{0.5em}{}
\titlespacing*{\section}{0pt}{14pt}{6pt}
\titlespacing*{\subsection}{0pt}{10pt}{4pt}

% ── Hyperlinks ────────────────────────────────────────────────────────────────
\usepackage[colorlinks=true, linkcolor=black, citecolor=black, urlcolor=blue]{hyperref}

% ── Miscellaneous ─────────────────────────────────────────────────────────────
\usepackage{enumitem}
\setlist[itemize]{noitemsep, topsep=3pt}
\usepackage{xcolor}
\usepackage{fancyhdr}

\pagestyle{fancy}
\fancyhf{}
\rhead{\small EE559 — Midway Progress Report}
\lhead{\small Om Suresh Prajapati}
\cfoot{\thepage}
\renewcommand{\headrulewidth}{0.4pt}

% ─────────────────────────────────────────────────────────────────────────────
\begin{document}

\begin{center}
    {\LARGE \textbf{Midway Progress Report}}\\[0.4em]
    {\large Adaptive Fraud Detection with Dynamic Thresholding\\
    and Cost-Sensitive Learning}\\[0.6em]
    \textbf{Om Suresh Prajapati} \quad | \quad USC ID: 1430823821\\[0.2em]
    \textit{EE559 (CSCI 559) — Machine Learning, Spring 2026}\\[0.2em]
    \textit{Instructor: Prof. Anand A. Joshi, University of Southern California}\\[0.2em]
    April 13, 2026
\end{center}

\vspace{0.3em}
\hrule
\vspace{0.6em}

% ─────────────────────────────────────────────────────────────────────────────
\section{Introduction and Motivation}

Credit card fraud causes billions of dollars in annual losses worldwide. Detecting fraud
is a supervised binary classification problem, but it presents a fundamental challenge:
the extreme imbalance between fraudulent and legitimate transactions. In the dataset used
for this project, only 0.17\% of transactions are fraudulent — a ratio of approximately
577 legitimate transactions for every 1 fraud case.

This class imbalance has two important consequences. First, it makes standard accuracy a
misleading metric: a classifier that predicts every transaction as legitimate achieves
99.83\% accuracy while being completely useless. Second, standard training algorithms tend
to ignore the minority class because it contributes so little to the overall loss. This
project addresses both consequences through appropriate evaluation metrics and
class-balancing strategies.

Beyond standard fraud classification, the key novelty of this project is the study of
\textit{dynamic decision thresholding}: rather than using a fixed threshold of 0.5 to
convert predicted probabilities to binary labels, we plan to investigate how the optimal
threshold changes when different costs are assigned to false positives (incorrectly flagging
a legitimate transaction) and false negatives (missing a fraudulent one). This
cost-sensitive analysis makes the system more practically useful.

% ─────────────────────────────────────────────────────────────────────────────
\section{Dataset Description and Exploratory Analysis}

\subsection{Dataset Details}

The dataset used is the \textbf{Kaggle Credit Card Fraud Detection Dataset}, a publicly
available benchmark for imbalanced binary classification \cite{kaggle_ccfd}.

\begin{itemize}
    \item \textbf{Task type:} Binary classification (fraud = 1, legitimate = 0)
    \item \textbf{Total samples:} 284,807 transactions
    \item \textbf{Features:} 30 input variables — 28 anonymized PCA components
          (\texttt{V1}--\texttt{V28}), plus \texttt{Time} and \texttt{Amount}
    \item \textbf{Class distribution:} 492 fraud cases (0.173\%), 284,315 legitimate (99.827\%)
    \item \textbf{Missing values:} None
    \item \textbf{Duplicate rows:} 1,081 exact duplicates (retained for now; to be investigated)
\end{itemize}

Features \texttt{V1}--\texttt{V28} are the result of a Principal Component Analysis (PCA)
transformation applied to protect cardholder privacy. They are already zero-centered and
decorrelated. Only \texttt{Time} (seconds since the first transaction in the dataset) and
\texttt{Amount} (transaction value in EUR) are in their original units.

\subsection{Key Findings from EDA}

\paragraph{Class imbalance severity.}
Figure~\ref{fig:class_dist} illustrates the 577:1 imbalance. Standard accuracy is
therefore discarded as a metric for the remainder of this project. Primary metrics are
\textbf{PR-AUC} (Precision-Recall Area Under Curve) and \textbf{ROC-AUC}.

\paragraph{Time and Amount features.}
Fraudulent transactions occur throughout the 48-hour observation window without a
clear temporal cluster, suggesting that time alone is a weak predictor of fraud.
Transaction amounts for fraud are notably smaller on average (mean: \$122 vs \$88 for
fraud), consistent with the common pattern of small test charges before larger fraudulent
withdrawals (Figure~\ref{fig:time_amount}).

\paragraph{Feature correlation with fraud.}
An absolute Pearson correlation analysis (Figure~\ref{fig:corr}) identified features
\texttt{V17}, \texttt{V14}, \texttt{V12}, \texttt{V10}, and \texttt{V11} as the most
linearly correlated with the fraud label. These features also show the clearest separation
between fraud and legitimate distributions in Figure~\ref{fig:feat_dist}, confirming
their discriminative value.

\begin{figure}[h]
    \centering
    \includegraphics[width=0.75\textwidth]{figures/fig_01_class_distribution.png}
    \caption{Severe class imbalance: 577 legitimate transactions for every 1 fraudulent.
             Standard accuracy is an uninformative metric under this distribution.}
    \label{fig:class_dist}
\end{figure}

\begin{figure}[h]
    \centering
    \includegraphics[width=0.90\textwidth]{figures/fig_02_time_amount.png}
    \caption{Left: Transaction time distribution by class. Right: Transaction amount
             distribution (log scale). Fraud amounts cluster at smaller values.}
    \label{fig:time_amount}
\end{figure}

% ─────────────────────────────────────────────────────────────────────────────
\section{Preprocessing Pipeline}

A leakage-free preprocessing pipeline is critical for valid model evaluation. The
following strict ordering was enforced:

\begin{enumerate}[noitemsep]
    \item \textbf{Split first:} Data was split into train (70\%), validation (15\%), and
          test (15\%) sets using \textit{stratified} sampling to preserve the 0.17\% fraud
          prevalence in all three splits. The test set is held out completely and will only
          be used for the final evaluation.
    \item \textbf{Fit transformers on training data only:} A \texttt{StandardScaler} was
          fitted exclusively on the training split for \texttt{Time} and \texttt{Amount}.
          The same fitted scaler was then applied to transform the validation and test sets.
    \item \textbf{No resampling for this report:} SMOTE and random under/over-sampling are
          planned for the final project phase. Applying them before splitting would
          constitute data leakage.
\end{enumerate}

\noindent\textbf{Why scale only \texttt{Time} and \texttt{Amount}?}
Features \texttt{V1}--\texttt{V28} are PCA outputs with zero mean and unit variance by
construction. Scaling them again would be redundant. Only the two original-scale features
require normalization before feeding into gradient- or distance-sensitive models
(Logistic Regression, SVM).

\noindent The resulting split sizes are shown in Table~\ref{tab:splits}.

\begin{table}[h]
    \centering
    \caption{Stratified train/validation/test split. Fraud prevalence is maintained
             at $\approx 0.17\%$ across all splits.}
    \label{tab:splits}
    \begin{tabular}{lccc}
        \toprule
        \textbf{Split} & \textbf{Samples} & \textbf{Fraud cases} & \textbf{Fraud \%} \\
        \midrule
        Train      & 199,020 & 345 & 0.173\% \\
        Validation & 42,722  & 74  & 0.173\% \\
        Test       & 43,065  & 73  & 0.170\% \\
        \midrule
        Total      & 284,807 & 492 & 0.173\% \\
        \bottomrule
    \end{tabular}
\end{table}

% ─────────────────────────────────────────────────────────────────────────────
\section{Baseline Model Results}

Two baseline classifiers were trained and evaluated on the \textbf{validation set}.
All reported metrics use a decision threshold of 0.50 for this midway report.

\subsection{Logistic Regression (L2 Regularization)}

Logistic Regression serves as the linear baseline. The model maps the input features
to a fraud probability via:
\[
P(y=1 \mid \mathbf{x}) = \sigma\!\left(\mathbf{w}^\top \mathbf{x} + b\right),
\quad
\mathcal{L} = \text{cross-entropy} + \frac{1}{2C}\|\mathbf{w}\|_2^2
\]
where $C = 1.0$ is the inverse regularization strength. The \texttt{class\_weight=`balanced'}
parameter adjusts sample weights inversely proportional to class frequencies, giving each
fraud sample $\approx 577\times$ more influence on the loss.

\subsection{Random Forest}

Random Forest is an ensemble of 100 decision trees, each trained on a bootstrap sample
with a random feature subset at each split. It can capture nonlinear feature interactions
that Logistic Regression cannot model. The same \texttt{class\_weight=`balanced'} setting
is applied at each tree node.

\subsection{Comparison}

Table~\ref{tab:results} summarizes both models on the validation set.
Figures~\ref{fig:lr_eval} and~\ref{fig:rf_eval} show the confusion matrices, ROC curves,
and Precision-Recall curves for each model.

\begin{table}[h]
    \centering
    \caption{Validation set performance (threshold = 0.50). ROC-AUC and PR-AUC are
             threshold-independent; Precision, Recall, and F1 are at threshold = 0.50.}
    \label{tab:results}
    \begin{tabular}{lccccc}
        \toprule
        \textbf{Model} & \textbf{ROC-AUC} & \textbf{PR-AUC} & \textbf{Fraud Recall} & \textbf{Fraud Precision} & \textbf{Fraud F1} \\
        \midrule
        Logistic Regression & -- & -- & -- & -- & -- \\
        Random Forest       & -- & -- & -- & -- & -- \\
        \bottomrule
        \multicolumn{6}{l}{\small\textit{Note: Fill in the values above after running the notebook.}} \\
    \end{tabular}
\end{table}

\begin{figure}[h]
    \centering
    \includegraphics[width=0.95\textwidth]{figures/fig_05_lr_evaluation.png}
    \caption{Logistic Regression: confusion matrix (left), ROC curve (center),
             and Precision-Recall curve (right) on the validation set.}
    \label{fig:lr_eval}
\end{figure}

\begin{figure}[h]
    \centering
    \includegraphics[width=0.95\textwidth]{figures/fig_06_rf_evaluation.png}
    \caption{Random Forest: confusion matrix (left), ROC curve (center),
             and Precision-Recall curve (right) on the validation set.}
    \label{fig:rf_eval}
\end{figure}

\paragraph{Key observation.}
At threshold 0.50, Logistic Regression tends to achieve higher Fraud Recall (catching more
actual fraud cases) while Random Forest achieves higher Fraud Precision and PR-AUC (fewer
false alarms, better overall precision-recall trade-off). Neither model is unambiguously
``better'' without a cost assumption — which is precisely the motivation for the
dynamic thresholding analysis planned for the final submission.

% ─────────────────────────────────────────────────────────────────────────────
\section{Challenges and Planned Next Steps}

\subsection{Challenges Encountered}

\begin{itemize}
    \item \textbf{Metric selection:} Choosing the right evaluation metric was a non-trivial
          design decision. ROC-AUC can be deceptively high on imbalanced datasets because the
          large number of true negatives inflates the denominator of the False Positive Rate.
          PR-AUC is more informative in this regime and will remain our primary metric.

    \item \textbf{Threshold ambiguity:} At threshold 0.50, Logistic Regression flags many
          false positives due to the model being biased toward high recall. The optimal
          threshold is highly application-dependent (what is the cost of a missed fraud vs.
          a false alarm?). This will be addressed via cost-sensitive threshold tuning.

    \item \textbf{PCA opacity:} Features V1--V28 have no semantic meaning, making feature
          engineering and interpretation harder. We will rely on Random Forest feature
          importances and SHAP values (planned) for model explainability.
\end{itemize}

\subsection{Planned Work for Final Submission}

\begin{enumerate}[noitemsep]
    \item \textbf{Third model:} Gradient Boosted Trees (XGBoost or LightGBM) or Support
          Vector Machine as a third baseline for comparison.
    \item \textbf{Imbalance handling:} SMOTE over-sampling and random under-sampling,
          applied strictly to the training set to avoid leakage.
    \item \textbf{Hyperparameter tuning:} Grid search with cross-validation on
          regularization strength, tree depth, and number of estimators.
    \item \textbf{Dynamic thresholding:} Sweep the decision threshold on the validation
          set; report the optimal threshold under multiple cost-ratio assumptions
          (FP cost : FN cost = 1:10, 1:50, 1:100).
    \item \textbf{Test set evaluation:} Apply the best model and optimal threshold to the
          held-out test set exactly once.
    \item \textbf{Error analysis:} Examine false negatives (missed frauds) for patterns
          in Amount, Time, or specific V-features.
\end{enumerate}

% ─────────────────────────────────────────────────────────────────────────────
\begin{thebibliography}{9}

\bibitem{kaggle_ccfd}
    ULB Machine Learning Group,
    \textit{Credit Card Fraud Detection Dataset},
    Kaggle, 2018.
    \url{https://www.kaggle.com/datasets/mlg-ulb/creditcardfraud}

\bibitem{sklearn}
    F. Pedregosa et al.,
    \textit{Scikit-learn: Machine Learning in Python},
    Journal of Machine Learning Research, 12, 2825--2830, 2011.

\end{thebibliography}

\end{document}
